\SourceLeaf{11}
They are also the solutions of the congruence
\[
1+(x+iy)^{p+1}\equiv0\quad\text{for}\quad p\equiv3\quad(\mathrm{mod.}\ 4).
\]
\par\indent From one value $x+iy$ one obtains them all by substituting for $u$ all values\\
$0,1,2,3,\, .\, .\, .,p-1$, using the formula\HistoricalNote{[*)]}{[*)}{The following formulas appear alongside this derivation in the manuscript:]
\[
1-xx\equiv yy,\qquad x=1-uy,\qquad 2u-uuy=y
\]
\[
\frac{1-uu}{1+uu}=x,\qquad\frac{2u}{1+uu}=y,
\]
\noindent [which explain the role of the additional factor $\dfrac{uu-1+2iu}{uu+1}$.]}
\[
(x+iy)\left(\frac{uu-1+2iu}{uu+1}\right)\quad\text{or}\quad\frac{(x+iy)(u+i)}{u-i}.
\]
\par\indent For $p\equiv1\ (\mathrm{mod.}\ 4)=aa+bb$, the formula $\dfrac{b}{a}.\dfrac{u+i}{u-i}$ contains all values of $x+iy$, where the only excluded values of $u$ are $\dfrac{a}{b}$ and $\dfrac{b}{a}$.
